\documentclass[12pt]{article}
\usepackage[utf8]{inputenc}
\usepackage[T1]{fontenc}
\usepackage{hyperref}
\title{Yagaria Language Resources for K-12 Curriculum Development}
\author{Compiled by Claude Sonnet 5.5 (claude-sonnet-5-5) for Andr\'as Kornai}
\date{}
\begin{document}
\maketitle

\section{Introduction}
Yagaria (ISO 639-3: ygr) is a Trans-New Guinea language of the Goroka District, Eastern Highlands Province, Papua New Guinea. Sources found in search disagree on the number of speakers: Wikipedia cites 38,000 (2000), while Joshua Project gives 52,000. The survey brief's figure of about 100,000 could not be confirmed. Little material exists, and the search found no school curriculum materials in Yagaria. This report lists only what the search returned.

\section{Primary Educational Resources}
No primers, textbooks or school programmes in Yagaria were found in search results.

\section{Linguistic and Reference Materials}
\begin{itemize}
\item Wikipedia overview of the language, including its dialect names: \url{https://en.wikipedia.org/wiki/Yagaria_language}
\item WALS entry for Yagaria. It lists Renck's \emph{Yagaria Dictionary with English Index} (1977) and Renck's 1975 grammar. These are older print works, and their current availability was not checked: \url{https://wals.info/languoid/lect/wals_code_ygr}
\end{itemize}

\section{Religious and Community Resources}
\begin{itemize}
\item Joshua Project people-group profile. It reports Bible portions translated, with no complete Bible: \url{https://joshuaproject.net/people_groups/15938}
\end{itemize}

\section{Digital and Audio Resources}
\begin{itemize}
\item Global Recordings Network page for a Yagaria variety (Dagenava). It offers evangelistic audio and does not offer recordings of the language as a linguistic resource: \url{https://glorecnet-website.k.globalrecordings.net/en/language/18391}
\end{itemize}

\section{Curriculum Development Context}
Yagaria is described as stable, but no evidence was found that it is taught in schools. A curriculum would have to start from the Renck dictionary and grammar and the Bible portions. It would need new orthography-based primers developed with community speakers.
\end{document}